\section{KV-Cache}

\subsection{¿Por qué se necesita el KV-Cache?}

\begin{importantbox}{Motivación central del KV-Cache}
En la generación autorregresiva, cada token nuevo necesita calcular la atención con \textbf{todos los tokens históricos}. Si en cada paso se recalculan todos los $K, V$, el costo computacional es de $O(M \cdot d_k)$ (donde $M$ crece continuamente). El KV-Cache almacena los vectores $K, V$ históricos en la memoria de la GPU, evitando el cálculo repetido.
\end{importantbox}

\subsection{Consumo de memoria del KV-Cache}

Tamaño del KV-Cache por token en cada capa:
$$
\text{Per-token KV} = 2 \times n_{\text{kv\_heads}} \times d_h \times \text{bytes}
$$

KV-Cache total:
$$
\text{Total KV} = L \times M \times 2 \times n_{\text{kv\_heads}} \times d_h \times \text{bytes}
$$

donde $L$ es el número de capas y $M$ es la longitud de la secuencia. Por esto GQA (reducir el número de KV heads) es tan importante para la eficiencia de la inferencia.

\subsection{Resumen de la sección}

El KV-Cache es la optimización central de la inferencia autorregresiva, pues evita recalcular $K, V$. Su consumo de memoria crece linealmente con la longitud de la secuencia y es el principal cuello de botella de la inferencia con contexto largo.
